\section{Discussion and conclusion}
The final experiment tests decision value directly rather than requiring
a new or globally more accurate credit predictor. Nevertheless, removing
the prediction-accuracy requirement does not remove the need for useful
decision information. In this experiment, fewer paid pairs and exact
request reuse do not yield the required utility improvement over random,
stale credit, and free controls.

The negative result is compatible with several mechanisms: sparse gains,
noisy small-sample estimates, insufficient adaptation of old counts,
misspecified marginal outcome models, or a service horizon too short to
amortize feedback. The present comparisons do not uniquely separate these
explanations. Retuning the horizon, weight, or prior using the revealed
service outcomes would create a new exploratory result, not repair the
failed preregistered candidate.

Across the five formal routes, evidence for conditional credit changed
more readily than evidence for beneficial paid maintenance. This suggests
a useful evaluation discipline: compare every acquisition method against
the same predictor, inexpensive structural rules, and a service-quality
baseline, while charging initialization and feedback. Such a discipline
can reject an attractive mechanism claim without rejecting the broader
research question.

The cycle used 76.476912 GPU residence hours of a 160-hour allowance and
exhausted its five-formal-candidate limit. The latter, not the final
candidate failure alone, is the project stopping condition. No sixth
candidate, fresh positive-method confirmation, or second-task validation
was conducted. We release an evidence-bounded negative draft and audit
record, with no claim of submission readiness or general effectiveness.
